\section{결론}
\label{sec:conclusion}

\begin{revision}

본 논문은 연속 CoC의 의무와 해제 조건을 후속 증거에 연결하고, 사건 순서에서 모순과 근거 부족을 구분하여 최초 문제 사건ㆍ의무ㆍ사유를 추적하는 프레임워크의 적용 가능성을 다루었다. 새 행동의 근거 B와 앞선 의무 A의 해제 증거를 구분하고, A 미해제 확인과 해제 여부 미상을 서로 다른 결과로 기록한다.

144개 작성 시나리오의 최초 시험에서 기존 구조화 경로는 96/144, 최초 위반은 54/72였다. 정상 36건 중 30건은 미상으로 남았고 자동 텍스트 경로는 전 사례 미상이었다. 단순 FSM은 144/144였다. 최초 실패 48건을 분석한 후 별도 조건별 UPPAAL 모델은 같은 사례에서 판정 144/144, 최초 위반 72/72, 문제 위치ㆍ의무ㆍ사유 각 108/108을 보였다. 이는 개발 후 동일 사례 재시험이며 독립 일반화 성능이 아니다.

기존 Python--UPPAAL의 판정ㆍ오류 종류ㆍ최초 위치는 144건에서 일치했다. 반례 접두 경로의 미상 사유 차이 24건은 종료 질의로 진단했으며 원 파서는 변경하지 않았다. 자연 검토는 309인접쌍 목록 중 별도 선정한 27장면으로, 텍스트 $\kappa=0.027301$과 합의 사유 기록 부재 때문에 자연 오류 정답으로 쓰지 않았다. 기존 P4 10쌍은 주석 오류와 분리한 처리 회귀시험이고 보조 반응 분석도 차량 안전 성능으로 해석하지 않는다.

작성된 미해제 충돌 사례에서는 이전 의무의 보존이 해당 사건 관계를 판정하는 데 필요했다. 다만 이 관찰을 모든 CoC에 대한 필요성ㆍ효과나 UPPAAL의 정확도 우위로 일반화하지 않는다. 현재 결과는 명시된 프레임의 시간ㆍ조건 관계를 실행하고 검토 근거를 추적하는 방법 실증이다. 자동 의미 추출, 개발 미사용 독립 시험, 사람 판정 신뢰도, 자연 오류 성능과 물리적 안전은 후속 검증으로 남는다.
\end{revision}
